\scanpage{34}
\mspage{14}

corps, éventuellement fini). On espère que Raynaud tirera au clair cette théorie, et son analogue local --- ce qui exige d'ailleurs de surmonter des obstacles techniques considérables.

Mais en dim $n\geq 2$, c'est sans espoir, car Artin a remarqué que tous les \deleted{\unsure{$p$-groupes ét.}} groupes de coh. de $X$, à coeff. dans des $p$-groupes constants, s'annulent en dim $>n+1$. Donc \uline{en mettant les choses au mieux, on arrivera à construire des ``bons'' invariants cohomologiques par cette voie jusqu'en dim} $\dim X+1$ (et non $2\dim X$, qui est la bonne). Il faudra donc (si $X$ projective lisse) \inserted{\unsure{pour le moins}} compléter le tableau \unsure{par} heuristique par dualité de Poincaré \ldots\ On voit bien \uline{qu'il faudra}

\margin{33}
